\section{Why keep the 1D baseline and add a 2D phase portrait task?}

The scalar control-affine environment remains useful as the minimal benchmark for online neural control with self-talk. It gives a clean setting where the controller gradient path is transparent and the role of the local loss
\[
\ell_t = x_t^2 + \rho u_t^2
\]
is easy to interpret.

At the same time, the scalar system is too limited to express genuine phase-space geometry. In one dimension the state can only move left or right, so there is no notion of circulation, trajectory bending, or phase portrait structure. To study how a controller can exploit geometry with small control effort, we add a fixed planar environment.

\section{Planar double-well environment}

The first 2D benchmark in \texttt{exp0526\_onlineNeuralControl} is the planar double-well system
\[
\dot x_1 = x_2,\qquad
\dot x_2 = \alpha x_1 - \beta x_1^3 - \gamma x_2 + c u.
\]

Here:
\begin{itemize}
    \item $x_1$ behaves like position,
    \item $x_2$ behaves like velocity,
    \item $u$ is a scalar control input,
    \item $-\gamma x_2$ provides damping,
    \item $\alpha x_1 - \beta x_1^3$ creates a double-well structure.
\end{itemize}

With the default choice
\[
\alpha = 1,\qquad \beta = 1,\qquad \gamma = 0.4,\qquad c = 1,
\]
the uncontrolled system has equilibria at
\[
(-1,0),\qquad (0,0),\qquad (1,0).
\]

The outer points $(-1,0)$ and $(1,0)$ are stable, while the center point $(0,0)$ is unstable. This gives the geometry we want:
\begin{itemize}
    \item the origin is not the natural resting point,
    \item there are stable nonzero operating points,
    \item small control can influence which basin the trajectory enters.
\end{itemize}

\section{Observation choice for the first 2D version}

The first 2D version uses full-state observation. The controller input is
\[
[x_1,\; x_2,\; 1],
\]
namely the full state plus the constant pulse channel already used throughout the project.

This choice is deliberate. The goal of the first 2D version is not partial observation or state estimation. The goal is to expose the controller to a system with richer geometry while keeping the observation path simple enough that rollout behavior and phase-space movement remain easy to read.

\section{Why phase portrait only in eval plots?}

The phase portrait is an analysis tool, not a training diagnostic. Training curves and training timelines are meant to answer:
\begin{itemize}
    \item how the loss evolves,
    \item how controls change after updates,
    \item how rollout windows are reshaped through optimization.
\end{itemize}

The phase portrait answers a different question:
\begin{itemize}
    \item where trajectories move in state space,
    \item which basin they approach,
    \item how blink or stutter perturbations deflect the path.
\end{itemize}

For that reason the phase portrait is added only to \texttt{eval\_rollout\_diagnostics.png}. It belongs to the evaluation layer, where we compare full, blink, and stutter rollouts against the vector field and equilibrium structure of the environment.
